\clearpage
\fancyhead[L]{\small BEMERKUNGEN VON SCHERING.}
\hypertarget{native-659-lower}{}
\begin{center}\textsc{Bemerkungen.}\end{center}

Dieser VI. Band von \textsc{Gauss} Werken enthält die von \textsc{Gauss} der Oeffentlichkeit übergebenen Abhandlungen, Aufsätze, Beobachtungen und Recensionen, in so fern sie den Astronomischen Wissenschaften angehören. Der Band ist dadurch schon stärker geworden als jeder der übrigen Bände dieses Werkes. Die im Nachlasse befindlichen dieses Gebiet betreffenden Handschriften, insbesondere eine Theorie der Planetarischen Störungen und die Pallastafeln konnten also nicht wol noch angeschlossen, sondern müssen von diesem Bande getrennt gedruckt werden.

Da die Briefe an \textsc{Zach} und \textsc{Lindenau} nicht mehr vorhanden sind, so habe ich aus den Zusammenstellungen, welche die Herausgeber der Monatlichen Correspondenz für Erd- und Himmelskunde von den Nachrichten über die kleinen Planeten veröffentlichten, diejenigen Stellen hier abdrucken lassen, welche \textsc{Gauss} betreffen und grossen Theils aus seinen Briefen entnommen waren, ausserdem nur noch solche Stellen, welche zum Verständniss der ersteren nothwendig erschienen. Die Ephemeriden habe ich nur in einzelnen wenigen Beispielen wiedergegeben. Bei den Vergleichungen von Elementen mit Beobachtungen habe ich diese letzteren beigefügt, wenn sie an anderen Orten in der Monatlichen Correspondenz sich finden. Für die aus dieser Zeitschrift entnommenen Beurtheilungen nichteigner Schriften ist der Name des Berichterstatters durch Briefe von \textsc{Lindenau} an \textsc{Gauss} bestätigt, für die aus den Göttinger gelehrten Anzeigen entnommenen durch die Rechnungsbücher der Königlichen Gesellschaft der Wissenschaften.

Göttingen 1874 November. \hfill \textsc{Ernst Schering}.
